\documentclass[12pt]{article}

% Formato de artigo da SBC fornecido em resources/template-latex.
\usepackage{sbc-template}
\usepackage[utf8]{inputenc}
\usepackage[brazil]{babel}
\usepackage{graphicx,url,amsmath}
\sloppy

\title{Comparação de genes candidatos à ortologia em quatro linhagens de \textit{Maribacter}}
\author{Autores do projeto\inst{1}}
\address{Instituição a informar}

\begin{document}
\maketitle

\begin{abstract}
  This report compares the \textit{recA}, \textit{trhO}, and \textit{rpoB} coding
  sequences and proteins in four \textit{Maribacter} lineages. We use global
  pairwise protein alignment, a progressive multiple alignment, codon
  projection, and parallel processing of six unique pairs per gene. Protein
  identity ranges from 86.30\% to 98.11\%. The evidence supports the genes as
  orthology candidates, but does not establish their evolutionary history.
\end{abstract}

\begin{resumo}
  Este trabalho compara as sequências codificadoras e as proteínas de
  \textit{recA}, \textit{trhO} e \textit{rpoB} em quatro linhagens de
  \textit{Maribacter}. São usados alinhamento global entre pares, alinhamento
  múltiplo progressivo, projeção de códons e processamento paralelo dos seis
  pares únicos por gene. A identidade proteica observada varia de 86,30\% a
  98,11\%. Os dados apoiam a hipótese de ortologia dos genes, sem demonstrar
  definitivamente sua história evolutiva.
\end{resumo}

\section{Introdução e objetivo}

O objetivo é comparar genes anotados em quatro linhagens do gênero
\textit{Maribacter} e descrever as diferenças observadas no DNA codificador e
nas proteínas correspondentes. O painel contém \textit{recA}, associado à
recombinação, \textit{trhO}, anotado como hidroxilase de uridina do tRNA, e
\textit{rpoB}, que codifica a subunidade beta da RNA polimerase. Uma CDS
(\textit{coding sequence}) é a região de DNA que codifica um produto proteico;
neste trabalho, cada CDS inclui o códon de parada terminal.

\subsection{Premissa e alcance da ortologia}

Adotou-se inicialmente, \textbf{sem uma reconstrução filogenética própria e,
portanto, com fundamento ainda limitado}, a premissa de que as linhagens
selecionadas compartilham um ancestral comum e seriam candidatos à ortologia. Essa foi uma hipótese de trabalho para escolher as sequências, não um resultado demonstrado. Em sentido estrito, \textbf{ortologia é uma relação entre genes de linhagens distintas que derivam
de um gene ancestral separado por especiação}; não é uma propriedade atribuída
a genomas inteiros. Genes semelhantes também podem ser parálogos, por exemplo
após duplicação \cite{tatusov1997,ncbiorthology}.

A avaliação posterior encontrou uma cópia anotada de cada gene por montagem,
CDSs com tradução íntegra e idêntica à proteína RefSeq correspondente, alta
identidade entre proteínas e vizinhança gênica conservada. Em uma triagem
heurística por segmentos idênticos de quatro aminoácidos, a proteína esperada
ficou em primeiro lugar nas 12 buscas direcionais entre proteomas para cada
gene. \textit{recA} e \textit{trhO} são vizinhos nas quatro montagens; a
vizinhança de \textit{rpoB} inclui \textit{rpoC}, \textit{rplL}, \textit{rplJ}
e \textit{rplA}. Essas observações sustentam os três grupos como candidatos
fortes à ortologia, mas a triagem não equivale a uma busca recíproca formal
com BLAST nem a uma comparação entre árvores de genes e de espécies. Duplicações
antigas seguidas de perdas ou transferência horizontal ainda são explicações
possíveis \cite{tatusov1997,ncbiorthology}.

\section{Dados e ambiente de execução}

Os genomas foram obtidos do NCBI. A Tabela~\ref{tab:montagens} identifica as
quatro montagens. Os arquivos dos genomas completos estão em:
\begin{itemize}
\item \path{resources/genomes/GCF_027474845.2_ASM2747484v2_genomic.fna};
\item \path{resources/genomes/GCF_002269385.1_ASM226938v1_genomic.fna};
\item \path{resources/genomes/GCF_001865655.1_ASM186565v1_genomic.fna};
\item \path{resources/genomes/GCF_001913155.1_ASM191315v1_genomic.fna}.
\end{itemize}
Para cada gene e montagem, o projeto extraiu uma CDS e a respectiva proteína,
usando as anotações RefSeq e os acessos registrados em
\path{resources/proteins.csv}. As 12 CDSs começam com \texttt{ATG}, terminam
em códon de parada, não contêm parada interna e traduzem exatamente para as
proteínas indicadas no RefSeq. A análise lê os FASTA de CDS (\texttt{.fna}) e
proteínas (\texttt{.faa}) em \path{resources/recA/}, \path{resources/trhO/}
ou \path{resources/rpoB/}, conforme o gene selecionado. Por exemplo, para
\textit{recA} da montagem \texttt{GCF\_001865655.1}, os arquivos de entrada
são \path{resources/recA/GCF_001865655.1_recA.fna} e
\path{resources/recA/GCF_001865655.1_recA.faa}. Os genomas completos
\textbf{não são alinhados integralmente}.

\begin{table}[ht]
\centering\small
\caption{Linhagens analisadas e montagens RefSeq.}
\label{tab:montagens}
\begin{tabular}{lll}
\hline
Linhagem usada no projeto & Montagem & Isolado \\
\hline
\textit{M. aestuarii} & \texttt{GCF\_027474845.2} & JCM18631 \\
\textit{M. cobaltidurans} & \texttt{GCF\_002269385.1} & B1 \\
\textit{M. dokdonensis} & \texttt{GCF\_001865655.1} & 1\_2014MBL\_MicDiv \\
\textit{M. hydrothermalis} & \texttt{GCF\_001913155.1} & T28 \\
\hline
\end{tabular}
\end{table}

A denominação \textit{M. dokdonensis} para \texttt{GCF\_001865655.1} é a
adotada no projeto e foi confirmada pelo responsável. O registro NCBI
\texttt{NZ\_CP011318.1} apresenta a mesma linhagem como \textit{Maribacter}
sp. 1\_2014MBL\_MicDiv \cite{ncbiassembly}.

A implementação principal usa \textbf{C++17}, biblioteca padrão de C++ e
\texttt{std::thread}; o \texttt{Makefile} compila com \texttt{g++},
\texttt{-O2} e \texttt{-pthread}. Scripts auxiliares de extração e conferência
usam Python 3. Como referência de execução e inspeção deste projeto, o
ambiente é Ubuntu 22.04.5 LTS em WSL2, arquitetura x86\_64, \texttt{g++}
11.4.0 e Python 3.10.12. Esses dados descrevem o ambiente consultado; os
resultados numéricos deste artigo foram lidos dos arquivos gerados em
\path{results/recA/}, \path{results/trhO/} e \path{results/rpoB/}.

Na raiz do repositório, \texttt{make test} executa as verificações e
\texttt{make} compila \texttt{build/gene-comparison}. Exemplos:
\begin{quote}\ttfamily\small
./build/gene-comparison --protein recA --benchmark\\
./build/gene-comparison --protein trhO --mode cds\\
./build/gene-comparison --protein rpoB --threads 1
\end{quote}
\texttt{--protein} aceita \texttt{recA}, \texttt{trhO} ou \texttt{rpoB}, com
\texttt{recA} como padrão. \texttt{--mode protein|cds} seleciona a unidade da
tabela impressa no terminal; os arquivos persistidos sempre contêm as duas.
\texttt{--input} e \texttt{--output-dir} mudam os diretórios padrão.

\section{Implementação dos alinhamentos}

\subsection{Alinhamento global entre pares}

As proteínas completas de cada gene, não os genomas completos, são alinhadas
pelo método global de Needleman e Wunsch \cite{needleman1970}. Para sequências
$x$ e $y$, a matriz de programação dinâmica $F$ representa a melhor pontuação
para os prefixos $x_{1..i}$ e $y_{1..j}$:
\begin{equation}
F(i,j)=\max\left\{\begin{array}{l}
F(i-1,j-1)+s(x_i,y_j),\\
F(i-1,j)-3,\\
F(i,j-1)-3,
\end{array}\right.
\qquad
s(a,b)=\begin{cases}2,&a=b,\\-1,&a\ne b.\end{cases}
\label{eq:nw}
\end{equation}
As bordas recebem penalidade de $-3$ por posição de lacuna; em empates, a
reconstrução prioriza a diagonal. As pontuações $+2$, $-1$ e $-3$ são escolhas
do projeto, e a lacuna usa penalidade linear. O cálculo guarda apenas duas
linhas de pontuações, além de um byte por célula para reconstruir o caminho.
Para comprimentos $m$ e $n$, o tempo é $O(mn)$ e a memória é $O(mn)$ por causa
da matriz de direções. A identidade relatada é a proporção de caracteres
idênticos nas colunas alinhadas, incluindo colunas com lacuna no denominador;
não é uma distância evolutiva corrigida \cite{yang2006}.

\subsection{Alinhamento múltiplo e códons}

Após os seis alinhamentos entre pares, o programa constrói separadamente um
alinhamento múltiplo \textbf{progressivo} das quatro proteínas, acrescentando
as sequências em ordem lexicográfica dos acessos de montagem. A ideia geral é
inspirada em Feng e Doolittle \cite{feng1987}, mas esta implementação não
constrói árvore guia; portanto, o resultado pode depender da ordem de entrada.
As colunas de aminoácidos orientam a projeção dos códons correspondentes das
CDSs, preservando trincas e o quadro de leitura \cite{suyama2006}. O códon de
parada é acrescentado ao fim do alinhamento de DNA e não integra o alinhamento
proteico. Para cada par, o programa conta bases diferentes, códons diferentes
que preservam o aminoácido e códons com aminoácidos diferentes.

\subsection{Paralelismo e otimizações}

Com $n=4$ sequências há $\binom{4}{2}=n(n-1)/2=6$ pares únicos: cada par
$(i,j)$ é calculado somente quando $i<j$. Não se repete a comparação inversa
$(j,i)$. Um contador atômico distribui esses seis trabalhos entre as threads;
o limite é o menor entre o valor de \texttt{--threads}, a concorrência de
hardware informada pelo sistema e o número de pares. Cada tarefa grava em
uma posição exclusiva do vetor de resultados, preservando a ordem de saída.
As sequências são carregadas uma vez e reutilizadas. O alinhamento múltiplo,
a projeção e a escrita dos relatórios ocorrem após a etapa paralela.
\texttt{--benchmark} repete apenas o alinhamento dos pares com uma thread e
confere igualdade dos resultados; os tempos \texttt{tempo\_pares\_paralelo\_ms}
e \texttt{tempo\_pares\_sequencial\_ms} abrangem somente essa etapa, não o fluxo
inteiro. Não se afirma aceleração sem medição controlada.

\section{Resultados}

\subsection{Visão geral}

A Tabela~\ref{tab:resumo} resume o alinhamento múltiplo. ``Conservadas''
significa colunas com a mesma letra nas quatro sequências; ``substituídas''
significa colunas sem lacuna e com pelo menos duas letras diferentes; ``com
lacuna'' significa colunas com ao menos um hífen. As três classes esgotam o
total de colunas em cada tipo de sequência.

\begin{table}[ht]
\centering\small
\caption{Resumo dos alinhamentos múltiplos de quatro sequências por gene.}
\label{tab:resumo}
\begin{tabular}{lrrrrrrrr}
\hline
& \multicolumn{4}{c}{Proteína (aa)} & \multicolumn{4}{c}{CDS (bases)} \\
Gene & Total & Cons. & Subst. & Lac. & Total & Cons. & Subst. & Lac. \\
\hline
\textit{recA} & 335 & 312 & 23 & 0 & 1008 & 693 & 315 & 0 \\
\textit{trhO} & 343 & 282 & 61 & 0 & 1032 & 651 & 381 & 0 \\
\textit{rpoB} & 1269 & 1202 & 67 & 0 & 3810 & 2815 & 995 & 0 \\
\hline
\end{tabular}
\end{table}

\textit{trhO} tem a maior fração de colunas proteicas com substituição
($61/343=17{,}78\%$), seguida de \textit{recA} ($23/335=6{,}87\%$) e
\textit{rpoB} ($67/1269=5{,}28\%$). Nenhum dos três alinhamentos tem lacunas
nesse painel. As diferenças no DNA são mais numerosas porque códons distintos
podem codificar o mesmo aminoácido. Essas são observações de sequências atuais,
sem atribuição de direção histórica ou efeito funcional.

\subsection{Identidades entre os seis pares}

Na Tabela~\ref{tab:pares}, os rótulos abreviados são D = linhagem
1\_2014MBL\_MicDiv, H = \textit{M. hydrothermalis}, C =
\textit{M. cobaltidurans} e A = \textit{M. aestuarii}. Em cada célula
``iguais/total (\%)'', o numerador conta colunas idênticas entre as duas
sequências, o denominador conta todas as colunas alinhadas e a porcentagem é
$100\times\mathrm{iguais}/\mathrm{total}$, arredondada a duas casas.

\begin{table}[ht]
\centering\small
\caption{Identidade observada por gene e par único, nas proteínas e CDSs.}
\label{tab:pares}
\begin{tabular}{llrr}
\hline
Gene & Par & Proteína: iguais/total (\%) & CDS: iguais/total (\%) \\
\hline
\textit{recA} & D--H & 328/335 (97,91) & 868/1008 (86,11) \\
 & D--C & 320/335 (95,52) & 803/1008 (79,66) \\
 & D--A & 322/335 (96,12) & 795/1008 (78,87) \\
 & H--C & 320/335 (95,52) & 803/1008 (79,66) \\
 & H--A & 323/335 (96,42) & 809/1008 (80,26) \\
 & C--A & 318/335 (94,93) & 786/1008 (77,98) \\
\hline
\textit{trhO} & D--H & 321/343 (93,59) & 855/1032 (82,85) \\
 & D--C & 305/343 (88,92) & 789/1032 (76,45) \\
 & D--A & 300/343 (87,46) & 787/1032 (76,26) \\
 & H--C & 296/343 (86,30) & 763/1032 (73,93) \\
 & H--A & 299/343 (87,17) & 799/1032 (77,42) \\
 & C--A & 299/343 (87,17) & 782/1032 (75,78) \\
\hline
\textit{rpoB} & D--H & 1245/1269 (98,11) & 3357/3810 (88,11) \\
 & D--C & 1225/1269 (96,53) & 3172/3810 (83,25) \\
 & D--A & 1224/1269 (96,45) & 3213/3810 (84,33) \\
 & H--C & 1224/1269 (96,45) & 3155/3810 (82,81) \\
 & H--A & 1224/1269 (96,45) & 3175/3810 (83,33) \\
 & C--A & 1234/1269 (97,24) & 3191/3810 (83,75) \\
\hline
\end{tabular}
\end{table}

Em todos os pares, a porcentagem de identidade da proteína supera a da CDS.
A comparação \textit{recA} D--A ilustra a distinção: há 213 bases diferentes
em códons codificadores alinhados, 173 posições de códons diferentes sem troca
de aminoácido e 13 posições com troca de aminoácido. Os números usam unidades
distintas: uma posição de códon pode conter duas ou três bases diferentes.
Por exemplo, \texttt{AGC} $\rightarrow$ \texttt{AGT} preserva serina;
\texttt{GAC} $\rightarrow$ \texttt{GAA} troca aspartato por glutamato. Em
\textit{trhO}, \texttt{AAA} $\rightarrow$ \texttt{ATA} troca lisina por
isoleucina. São diferenças entre sequências presentes, não mutações
historicamente localizadas \cite{yang2006,zhang1997}.

\subsection{Como ler cada arquivo e campo de saída}

O programa grava um diretório \texttt{results/<gene>/} para cada execução.
\texttt{protein\_alignment.faa} contém quatro linhas FASTA de proteínas
alinhadas; \texttt{cds\_alignment.fna} contém as quatro CDSs alinhadas por
códons. Cada cabeçalho FASTA identifica montagem, acesso de sequência e gene.
O caractere \texttt{-} representa lacuna. Como as CDSs incluem o códon de
parada, o alinhamento de DNA tem três colunas finais além de três colunas
por aminoácido.

\noindent\textbf{\texttt{pairs.tsv}.} Cada linha representa um dos seis pares
únicos do gene. \texttt{montagem\_1} e \texttt{montagem\_2} são os acessos
das montagens; \texttt{proteína\_1} e \texttt{proteína\_2} são os acessos
RefSeq das proteínas comparadas. \texttt{coincidências\_proteína} conta posições
com o mesmo aminoácido nos dois lados; \texttt{colunas\_proteína} é o número
total de colunas do alinhamento do par; \texttt{identidade\_proteína\_percentual}
é $100\times\texttt{coincidências\_proteína}/\texttt{colunas\_proteína}$;
\texttt{lacunas\_proteína} conta colunas com hífen em pelo menos um lado.
\texttt{coincidências\_cds}, \texttt{colunas\_cds},
\texttt{identidade\_cds\_percentual} e \texttt{lacunas\_cds} são as quatro medidas
análogas para bases, após projetar o alinhamento proteico sobre as CDSs. A
identidade da CDS inclui o códon de parada; a identidade proteica não.
Por exemplo, \textit{recA} D--H apresenta 328 aminoácidos idênticos em 335
colunas (97,91\%) e 868 bases idênticas em 1008 colunas (86,11\%), sem
lacunas.

\noindent\textbf{\texttt{codon\_effects.tsv}.} \texttt{montagem\_1} e
\texttt{montagem\_2} identificam o mesmo par de montagens. O campo
\texttt{bases\_diferentes} conta bases diferentes em posições codificadoras
nas quais os dois códons estão presentes; exclui o códon de parada e posições
com lacuna. \texttt{códons\_sinônimos} conta posições cujo códon difere,
mas cujo aminoácido é igual. \texttt{códons\_não\_sinônimos} conta posições
com aminoácidos diferentes. \texttt{códons\_com\_lacuna} conta posições de códon
com lacuna em uma das sequências. \texttt{parada\_diferente} vale
\texttt{sim} quando os tripletos terminais de parada diferem, mesmo que
ambos interrompam a tradução; \texttt{não} indica tripletos iguais. As
contagens sinônimas e não sinônimas são de \emph{posições de códon}, não de
bases ou de eventos evolutivos. Elas não são estimativas de $d_N$, $d_S$ ou
$d_N/d_S$.

\noindent\textbf{\texttt{summary.tsv}.} Há uma linha por gene.
\texttt{gene} identifica \textit{recA}, \textit{trhO} ou \textit{rpoB};
\texttt{sequências} é a quantidade de sequências no alinhamento múltiplo
(quatro neste estudo). \texttt{colunas\_proteína} e \texttt{colunas\_cds}
são os comprimentos do alinhamento múltiplo em aminoácidos e bases,
respectivamente. \texttt{proteína\_conservada} e \texttt{cds\_conservada}
contam colunas iguais nas quatro linhas; \texttt{proteína\_substituída} e
\texttt{cds\_substituída} contam colunas sem lacunas, mas com ao menos uma
diferença; \texttt{proteína\_com\_lacuna} e \texttt{cds\_com\_lacuna} contam
colunas com pelo menos uma lacuna. Essas medidas vêm do alinhamento
\emph{múltiplo}; não devem ser confundidas com as medidas por par de
\texttt{pairs.tsv}.

O terminal ainda mostra \texttt{gene}, \texttt{modo}, \texttt{sequências},
\texttt{pares\_únicos}, \texttt{threads} e \texttt{diretório\_resultados}: gene
selecionado, unidade exibida, número de sequências, pares calculados, threads
usadas e diretório dos arquivos. Na tabela do terminal, \texttt{acesso\_1}
e \texttt{acesso\_2} são os acessos FASTA correspondentes ao modo;
\texttt{coincidências}, \texttt{colunas}, \texttt{identidade\_percentual} e
\texttt{lacunas} seguem as mesmas definições por par. \texttt{pontuação} é a soma
das pontuações $+2/-1/-3$ no alinhamento exibido e depende do modo;
\texttt{bases\_diferentes}, \texttt{códons\_sinônimos},
\texttt{códons\_não\_sinônimos} e \texttt{códons\_com\_lacuna} mantêm as definições
acima. \texttt{colunas\_alinhamento\_múltiplo}, \texttt{conservadas},
\texttt{substituídas} e \texttt{com\_lacuna} no terminal resumem sempre o
alinhamento múltiplo de \emph{proteínas}. Se solicitado \texttt{--benchmark},
\texttt{resultados\_idênticos=sim} informa que os alinhamentos dos pares com
uma thread e com múltiplas threads coincidiram.

\section{Discussão, limitações e conclusão}

As altas identidades, a cópia anotada única, a tradução íntegra e a sintenia
conservada apoiam a seleção de \textit{recA}, \textit{trhO} e \textit{rpoB}
como candidatos à ortologia nas quatro linhagens. A premissa inicial de
ancestralidade comum motivou a comparação, mas não substitui uma análise
filogenética. A porcentagem de identidade é uma medida descritiva e pode
subestimar a quantidade de substituições históricas \cite{yang2006}.
Similarmente, diferenças sinônimas e não sinônimas não revelam por si só a
linhagem em que ocorreram, convergência, seleção ou mudança funcional;
inferências desse tipo exigiriam filogenia e reconstrução ancestral
\cite{zhang1997}.

O estudo entrega seis comparações globais por gene, alinhamentos múltiplos
das proteínas, alinhamentos de CDS orientados por códons e tabelas que
separam identidade de proteínas, identidade de DNA e efeitos observados
por códon. Uma continuação apropriada é comparar árvores gênicas com uma
árvore de espécies e ampliar a amostragem de linhagens antes de afirmar
ortologia definitiva ou interpretar mecanismos evolutivos.

\bibliographystyle{sbc}
\bibliography{referencias}
\end{document}
